A pollen grain jiggling in water, a stock price, the spectrum of a heavy nucleus and the zeros of the Riemann zeta function
all move to the same music: Brownian motion. This book builds it from Gaussian processes, studies the fine structure of its
paths, and develops the It\^o calculus that it requires. It then puts the motion to work: the heat equation and harmonic
functions, Girsanov's change of measure, stochastic differential equations from Ornstein--Uhlenbeck to Black--Scholes and
Dyson, L\'evy processes, random walks on networks as electrical circuits, and expander graphs with their own Riemann
hypothesis. The last part is about random matrices: Wigner's semicircle, the sine kernel of random unitary matrices, the
moments of their characteristic polynomials found by Keating and Snaith, and the log-correlated fields and multiplicative
chaos that describe how large these polynomials, and perhaps the zeta function, can get. Every result is proved, with two
marked exceptions at the research frontier. Each chapter opens with the story behind its ideas, ends with exercises and
complete solutions.
